\section{Gradient-Descent Attitude Estimation: The Madgwick Filter}
\label{sec:madgwick}

% ----------------------------------------------------------------
% Motivation, derivation, and comparison with the complementary filter.
% ----------------------------------------------------------------

The complementary filter presented in Section~\ref{sec:complementary_filter}
provides an effective and intuitive approach to sensor fusion.  However, its
scalar formulation operates on Euler angles independently, which introduces
two practical limitations: (i)~the relationship between body-frame angular
velocities and Euler angle rates is \emph{nonlinear and coupled} (requiring
explicit Euler rate equations), and (ii)~the representation becomes singular
near $\pm 90^\circ$ pitch (gimbal lock).  While the quaternion form of the
complementary filter (Section~\ref{sec:cf:quaternion}) addresses these issues
in principle via SLERP blending, an alternative family of algorithms performs
the fusion \emph{natively} in quaternion space using gradient descent on an
orientation error objective.  The most widely adopted of these is the filter
proposed by Madgwick et al.~(2011), which we describe in this section.

\subsection{Problem Statement}
\label{sec:madg:problem}

Given a stream of gyroscope readings
$\boldsymbol{\omega}_b \in \mathbb{R}^3$, accelerometer readings
$\mathbf{a}_b \in \mathbb{R}^3$, and (optionally) magnetometer readings
$\mathbf{m}_b \in \mathbb{R}^3$, we seek the unit quaternion
$\mathbf{q}(t)$ that represents the orientation of the sensor body frame
with respect to a fixed world frame at each time step.

The key idea is to formulate attitude estimation as an \emph{optimisation
problem}: find the quaternion that best aligns the sensor measurements with
known reference directions in the world frame (gravity and the geomagnetic
field).

\subsection{Gyroscope Prediction}
\label{sec:madg:gyro}

As in the complementary filter, the gyroscope provides a short-term
orientation prediction.  Given the current estimate $\mathbf{q}(t)$ and the
angular velocity $\boldsymbol{\omega}_b(t)$, the predicted quaternion is
obtained by integrating the quaternion kinematic equation
(Eq.~\ref{eq:quat_euler}):

\begin{equation}
  \mathbf{q}_{\omega}(t + \Delta t)
  = \mathbf{q}(t)
    + \frac{\Delta t}{2}\,\mathbf{q}(t) \otimes \boldsymbol{\omega}_q(t)
  \label{eq:madg_gyro}
\end{equation}

\noindent where $\boldsymbol{\omega}_q = (0,\, \omega_x,\, \omega_y,\,
\omega_z)$.  This step is identical to the gyroscope path in the
complementary filter and inherits the same drift characteristics.

\subsection{Accelerometer and Magnetometer Correction via Gradient Descent}
\label{sec:madg:gradient}

The correction step is where the Madgwick filter fundamentally differs from
the complementary filter.  Instead of extracting Euler angles from the
accelerometer and blending them with the gyroscope output, the Madgwick
filter defines an \emph{objective function} whose minimum corresponds to the
correct orientation.

\subsubsection{Objective Function (IMU: Accelerometer Only)}

Let $\hat{\mathbf{d}}_w = (0,\, 0,\, 1)^\top$ be the expected direction of
gravity in the world frame (normalised).  If the true orientation is
$\mathbf{q}$, then the accelerometer should measure:

\begin{equation}
  \hat{\mathbf{a}}_b = \mathbf{q}^* \otimes \hat{\mathbf{d}}_{w,q}
                       \otimes \mathbf{q}
  \label{eq:expected_accel}
\end{equation}

\noindent The error between the predicted and actual normalised accelerometer
reading $\hat{\mathbf{a}}_b^{\,\text{meas}}$ defines the objective function:

\begin{equation}
  f(\mathbf{q})
  = \mathbf{q}^* \otimes \hat{\mathbf{d}}_{w,q} \otimes \mathbf{q}
    - \hat{\mathbf{a}}_{b,q}^{\,\text{meas}}
  \label{eq:objective}
\end{equation}

\noindent where $f : \mathbb{R}^4 \to \mathbb{R}^3$ maps a quaternion to the
residual between the expected and observed gravity direction.  The optimal
orientation satisfies $f(\mathbf{q}) = \mathbf{0}$.

\subsubsection{Gradient Step}

Rather than solving the full nonlinear optimisation, Madgwick applies a
\emph{single gradient descent step} per time sample, which is sufficient
because the orientation changes slowly relative to the sampling rate:

\begin{equation}
  \mathbf{q}_{\nabla}(t)
  = \mathbf{q}(t)
    - \beta\,\frac{\nabla f(\mathbf{q})}{\|\nabla f(\mathbf{q})\|}
  \label{eq:gradient_step}
\end{equation}

\noindent where $\nabla f = J^\top\!(\mathbf{q})\, f(\mathbf{q})$ is
computed analytically from the Jacobian $J(\mathbf{q})$ of the objective
function with respect to the quaternion components.  The normalisation
$\|\nabla f\|$ ensures that the step magnitude depends only on $\beta$, not
on the error amplitude.

\subsubsection{Extension to MARG (with Magnetometer)}

When magnetometer data is available, the objective function is extended to
include the geomagnetic field direction.  Let $\hat{\mathbf{b}}_w$ be the
expected magnetic field direction in the world frame (computed from the
magnetometer measurement itself, projected to have zero East component).
The combined objective becomes:

\begin{equation}
  f_{\text{MARG}}(\mathbf{q})
  = \begin{bmatrix}
      f_g(\mathbf{q}) \\[2pt]
      f_b(\mathbf{q})
    \end{bmatrix}
  \label{eq:objective_marg}
\end{equation}

\noindent where $f_g$ is the gravity residual and $f_b$ is the magnetic field
residual.  The Jacobian becomes a $6 \times 4$ matrix, and the gradient step
proceeds analogously.  The magnetometer provides an absolute heading
reference that the accelerometer alone cannot observe (since gravity is
invariant under rotations about the vertical axis, as discussed in
Section~\ref{sec:proj:gravity}).

\subsection{Filter Fusion}
\label{sec:madg:fusion}

The final quaternion estimate is a weighted combination of the gyroscope
prediction and the gradient correction:

\begin{equation}
  \boxed{
    \mathbf{q}(t + \Delta t)
    = \gamma\,\mathbf{q}_{\omega}(t + \Delta t)
      + (1 - \gamma)\,\mathbf{q}_{\nabla}(t + \Delta t)
  }
  \label{eq:madg_fusion}
\end{equation}

\noindent followed by renormalisation $\mathbf{q} \leftarrow \mathbf{q} /
\|\mathbf{q}\|$.  In the original formulation, the gain parameter $\beta$
simultaneously controls the step size and the effective blending ratio,
playing a role analogous to $\alpha$ in the complementary filter.  A typical
value is $\beta = 0.041$ for MARG configurations.

The crucial difference is that \textbf{the entire fusion operates in
quaternion space}: there is no intermediate conversion to Euler angles, no
need for Euler rate equations, and no singularity at $\pm 90^\circ$ pitch.

\subsection{Comparison with the Complementary Filter}
\label{sec:madg:comparison}

Table~\ref{tab:cf_vs_madg} summarises the key differences between the two
approaches.

\begin{table}[t]
  \centering
  \caption{Complementary filter vs.\ Madgwick filter.}
  \label{tab:cf_vs_madg}
  \small
  \begin{tabular}{@{}lcc@{}}
    \toprule
    \textbf{Property}
      & \textbf{Complementary}
      & \textbf{Madgwick} \\
    \midrule
    Fusion domain
      & Euler angles
      & Quaternion \\
    Gimbal lock
      & Yes (scalar form)
      & No \\
    Euler rate equations
      & Required
      & Not needed \\
    Magnetometer fusion
      & Separate filter
      & Integrated \\
    Tuning parameters
      & $\alpha$ per axis
      & Single $\beta$ \\
    Computational cost
      & Very low
      & Low \\
    Gravity sign convention
      & Must be handled
      & Automatic \\
    \bottomrule
  \end{tabular}
\end{table}

\subsubsection{Euler Rate Coupling}

In the scalar complementary filter, the gyroscope rates
$(\omega_x, \omega_y, \omega_z)$ are used directly as the rates of change
of roll, pitch, and yaw.  This is only valid for small angles.  The correct
relationships are:

\begin{align}
  \dot{\phi}   &= \omega_x
    + \sin\phi\,\tan\theta\;\omega_y
    + \cos\phi\,\tan\theta\;\omega_z
  \label{eq:euler_rate_roll} \\[3pt]
  \dot{\theta} &= \cos\phi\;\omega_y - \sin\phi\;\omega_z
  \label{eq:euler_rate_pitch} \\[3pt]
  \dot{\psi}   &= \frac{\sin\phi}{\cos\theta}\;\omega_y
    + \frac{\cos\phi}{\cos\theta}\;\omega_z
  \label{eq:euler_rate_yaw}
\end{align}

\noindent At large inclinations (e.g.\ $\theta > 20^\circ$), ignoring this
coupling introduces significant cross-axis errors that corrupt the
orientation estimate and, consequently, the gravity removal step.  The
Madgwick filter avoids this entirely because quaternion integration
(Eq.~\ref{eq:madg_gyro}) naturally handles the nonlinear coupling
between axes without any explicit trigonometric corrections.

\subsubsection{Gravity Sign Convention}

A subtle but practically important difference concerns the sign of the
gravity vector.  The complementary filter extracts tilt angles from the
accelerometer and subtracts gravity as a scalar magnitude
$\mathbf{g}_w = (0, 0, \pm g)$, where the sign depends on the sensor
convention.  If the sensor reports $a_z \approx +9.8$~m/s$^2$ at rest
(screen up), the gravity vector must be subtracted as $(0, 0, +g)$; if it
reports $a_z \approx -9.8$~m/s$^2$ (screen down or alternate convention),
the sign must be flipped.

The Madgwick filter is immune to this issue: it works with the
\emph{normalised} accelerometer reading and optimises for the direction
of gravity, not its sign.  The gravity vector in the world frame is then
determined empirically from the calibration phase as:

\begin{equation}
  \mathbf{g}_w = \frac{1}{N_c} \sum_{n=1}^{N_c}
    \mathbf{a}_w[n]
  \label{eq:g_world_measured}
\end{equation}

\noindent where $\mathbf{a}_w[n] = \mathbf{q}[n] \otimes
\mathbf{a}_{b,q}[n] \otimes \mathbf{q}^*[n]$ is the accelerometer reading
rotated to the world frame using the Madgwick quaternion.  This approach is
robust to any sensor sign convention and eliminates a common source of
implementation errors.

\subsubsection{Effect on Position Estimation}

Because the Madgwick filter produces a more accurate orientation estimate
(especially at large inclinations), the gravity removal step is more
precise, which directly reduces the gravity leakage that dominates the
position drift.  In our experiments with a 39-second recording of repetitive
vertical hand movements (``subidas\_bajadas\_con\_pausa''), the Madgwick
pipeline produced:

\begin{itemize}
  \item A Z-axis range of $[0.00,\; 0.76]$~m, consistent with the actual
        hand movement amplitude.
  \item No negative Z values, confirming that the hand never descended below
        its starting position (as expected from the experimental protocol).
  \item Near-zero residual bias after gravity subtraction
        ($\sim 10^{-6}$~m/s$^2$), compared to biases of order $\pm 19$~m/s$^2$
        with the scalar complementary filter before sign correction.
\end{itemize}

\noindent Furthermore, the Madgwick pipeline does not require a high-pass
filter on the linear acceleration, which was previously necessary to
compensate for gravity leakage in the complementary filter approach.  The
removal of this filter is significant because the high-pass filter, while
effective at removing drift, also \emph{destroys sustained displacements}:
it forces the acceleration to have zero mean, which means any net position
change (e.g.\ moving up and staying up) is gradually pulled back toward the
origin.  This manifested as backward loops at rest positions and spurious
negative heights---artefacts that disappear entirely with the improved
orientation from the Madgwick filter.
